%%%%%%%%%%%%%%%%%%%%%%%%%%%%%%%%%%%%%%%%%
% Professional Formal Letter
% LaTeX Template
% Version 2.0 (12/2/17)
%
% This template originates from:
% http://www.LaTeXTemplates.com
%
% Authors:
% Brian Moses
% Vel (vel@LaTeXTemplates.com)
%
% License:
% CC BY-NC-SA 3.0 (http://creativecommons.org/licenses/by-nc-sa/3.0/)
%
%%%%%%%%%%%%%%%%%%%%%%%%%%%%%%%%%%%%%%%%%

%----------------------------------------------------------------------------------------
%	PACKAGES AND OTHER DOCUMENT CONFIGURATIONS
%----------------------------------------------------------------------------------------

\documentclass[11pt, a4paper]{letter} % Set the font size (10pt, 11pt and 12pt) and paper size (letterpaper, a4paper, etc)

%%%%%%%%%%%%%%%%%%%%%%%%%%%%%%%%%%%%%%%%%
% Professional Formal Letter
% Structure Specification File
% Version 1.0 (12/2/17)
%
% This file originates from:
% http://www.LaTeXTemplates.com
%
% Authors:
% Brian Moses
% Vel (vel@LaTeXTemplates.com)
%
% License:
% CC BY-NC-SA 3.0 (http://creativecommons.org/licenses/by-nc-sa/3.0/)
%
%%%%%%%%%%%%%%%%%%%%%%%%%%%%%%%%%%%%%%%%%

%----------------------------------------------------------------------------------------
%	PACKAGES AND OTHER DOCUMENT CONFIGURATIONS
%----------------------------------------------------------------------------------------

\usepackage{graphicx} % Required for including pictures

\usepackage[T1]{fontenc} % Output font encoding for international characters
\usepackage[utf8]{inputenc} % Required for inputting international characters
\usepackage{gfsdidot} % Use the GFS Didot font: http://www.tug.dk/FontCatalogue/gfsdidot/
\usepackage{microtype} % Improves typography

\pagestyle{empty} % Suppress headers and footers

\setlength\parindent{1cm} % Paragraph indentation

% Create a new command for the horizontal rule in the document which allows thickness specification
\makeatletter
\newcommand{\vhrulefill}[1]{\leavevmode\leaders\hrule\@height#1\hfill \kern\z@}
\makeatother

%----------------------------------------------------------------------------------------
%	DOCUMENT MARGINS
%----------------------------------------------------------------------------------------

\usepackage{geometry} % Required for adjusting page dimensions

\geometry{
	top=0.5in, % Top margin
	bottom=0.5in, % Bottom margin
	left=0.5in, % Left margin
	right=0.5in, % Right margin
	%showframe, % Uncomment to show how the type block is set on the page
}

%----------------------------------------------------------------------------------------
%	DEFINE CUSTOM COMMANDS
%----------------------------------------------------------------------------------------

\newcommand{\logo}[1]{\renewcommand{\logo}{#1}}

\newcommand{\Who}[1]{\renewcommand{\Who}{#1}}
\newcommand{\Title}[1]{\renewcommand{\Title}{#1}}

\newcommand{\headerlineone}[1]{\renewcommand{\headerlineone}{#1}}
\newcommand{\headerlinetwo}[1]{\renewcommand{\headerlinetwo}{#1}}

\newcommand{\authordetails}[1]{\renewcommand{\authordetails}{#1}}

%----------------------------------------------------------------------------------------
%	AUTHOR DETAILS STRUCTURE
%----------------------------------------------------------------------------------------

\newcommand{\authordetailsblock}{
	\hspace{\fill} % Move the author details to the far right
	\parbox[t]{0.48\textwidth}{ % Box holding the author details; width value specifies where it starts and ends, increase to move details left
		\footnotesize % Use a smaller font size for the details
		\Who\\ % Author name
		\textit{\authordetails} % The author details text, all italicised
	}
}

%----------------------------------------------------------------------------------------
%	HEADER STRUCTURE
%----------------------------------------------------------------------------------------

\address{
	\includegraphics[width=2in]{\logo} % Include the logo of author institution
% 	\hspace{0.\textwidth} % Position of the institution logo, increase to move left, decrease to move right
% 	\vskip -0.1\textheight~\\ % Position of the large header text in relation to the institution logo, increase to move down, decrease to move up
% 	\Large\hspace{0.2\textwidth}\headerlineone\hfill ~\\[0.006\textheight] % First line of institution name, adjust hspace if your logo is wide
% 	\hspace{0.2\textwidth}\headerlinetwo\hfill \normalsize % Second line of institution name, adjust hspace if your logo is wide
% 	\makebox[0ex][r]{\textbf{\Who\Title}}\hspace{0.01\textwidth} % Print author name and title with a little whitespace to the right
% 	~\\[-0.01\textheight] % Reduce the whitespace above the horizontal rule
% 	\hspace{0.2\textwidth}\vhrulefill{1pt} \\ % Horizontal rule, adjust hspace if your logo is wide and \vhrulefill for the thickness of the rule
	\authordetailsblock % Include the letter author's details on the right side of the page under the horizontal rule
	\hspace{-0.25\textwidth} % Horizontal position of the author details block, increase to move left, decrease to move right
	\vspace{-0.1\textheight} % Move the date and letter content up for a more compact look
}

%----------------------------------------------------------------------------------------
%	COMPOSE THE ENTIRE HEADER
%----------------------------------------------------------------------------------------

\renewcommand{\opening}[1]{
	{\centering\fromaddress\vspace{0.05\textheight} \\ % Print the header and from address here, add whitespace to move date down
	\hspace*{\longindentation}\today\hspace*{\fill}\par} % Print today's date, remove \today to not display it
	{\raggedright \toname \\ \toaddress \par} % Print the to name and address
	\vspace{1cm} % White space after the to address
	\noindent #1 % Print the opening line
}

%----------------------------------------------------------------------------------------
%	SIGNATURE STRUCTURE
%----------------------------------------------------------------------------------------

% \signature{\Who\Title} % The signature is a combination of the author's name and title

% \renewcommand{\closing}[1]{
% 	\vspace{2.5mm} % Some whitespace after the letter content and before the signature
% 	\noindent % Stop paragraph indentation
% 	\hspace*{\longindentation} % Move the signature right to the value of \longindentation
% 	\parbox{\indentedwidth}{
% 		\raggedright
% 		#1 % Print the signature text
% 		\vskip 1.65cm % Whitespace between the closing text and author's name for a physical signature
% 		\fromsig % Prints the value of \signature{}, i.e. author name and title
% 	}
% }
 % Include the file that specifies the document structure

%\longindentation=0pt % Un-commenting this line will push the closing "Sincerely," and date to the left of the page

%----------------------------------------------------------------------------------------
%	YOUR INFORMATION
%----------------------------------------------------------------------------------------

% \Who{Youngmin Park, PhD} % Your name

% \Title{} % Your title, leave blank for no title

\authordetails{
Youngmin Park, PhD\\
	Department of Otorhinolaryngology\\ % Your department/institution
	University of Pennsylvania\\ % Your address
	3400 Spruce St. 5 Ravdin\\
	Philadelphia, PA 19104\\ % Your city, zip code, country, etc
	Email: ympark1988@gmail.com\\ % Your email address
% 	Phone: (412) 805-0283\\ % Your phone number
	URL: youngmp.github.io % Your URL
}

%----------------------------------------------------------------------------------------
%	HEADER CONTENTS
%----------------------------------------------------------------------------------------

\logo{PerelmanLogoBig.jpg} % Logo filename, your logo should have square dimensions (i.e. roughly the same width and height), if it does not, you will need to adjust spacing within the HEADER STRUCTURE block in structure.tex (read the comments carefully!)

% \headerlineone{UNIVERSITY} % Top header line, leave blank if you only want the bottom line

% \headerlinetwo{OF PENNSYLVANIA} % Bottom header line

%----------------------------------------------------------------------------------------

\begin{document}

%----------------------------------------------------------------------------------------
%	TO ADDRESS
%----------------------------------------------------------------------------------------

\begin{letter}{}
% {
% 	Youngmin Park, PhD, Mathematics\\
% 	Department of Otorhinolaryngology\\
% 	University of Pennsylvania\\
% 	Philadelphia, PA 19104
% }

%----------------------------------------------------------------------------------------
%	LETTER CONTENT
%----------------------------------------------------------------------------------------

\opening{Dear Sir or Madam,}

I am applying for the postdoctoral position in the Program in Applied and Computational Mathematics (PACM). I completed a PhD in mathematics with advisor G. Bard Ermentrout, and I am now a postdoc at the University of Pennsylvania under a one-year contract. Under the guidance of advisor Maria N. Geffen, we are modeling the roles of inhibitory interneurons on complex auditory processing.

Within PACM, I am interested in working with Naomi Leonard, who has substantial works at the interface of dynamical systems and biology. Several of her publications, including ``Mixed mode oscillations and phase locking in coupled FitzHugh-Nagumo model neurons'' (2018), align with my goal of understanding the role of oscillations in biology, and determining general conditions for synchrony and asynchrony. This particular paper is of great personal interest, as it includes the exploration of fast-slow systems and the existence of mixed-mode oscillations and canards.

As proof of my qualification to carry out successful collaborations, I reference my first-author publications in mathematical neuroscience while working with my dissertation advisor, Bard Ermentrout. In our first paper, we explore the effects of a neurotransmitter with slowly varying concentration on the synchrony of weakly coupled neurons through a phase reduction. Our second paper extends this phase reduction to account for internally-generated slow processes, and naturally includes cases with large phase drifts between populations. Our third paper introduces a novel dimension-reduction of a non-local equation of neural activity, which allows for a rigorous and thorough classification of the many different bifurcations of spatio-temporal solutions of the system. In addition to these primary topics, we published a book chapter on the consequences of the shape of the phase response curve on the synchronization properties of inhibitory and excitatory neurons. The sum of this research resulted in winning the prestigious Andrew Mellon Predoctoral Fellowship, which is awarded to doctoral students at the University of Pittsburgh of exceptional promise and ability.

I remark that the bulk of this work was completed as a doctoral student while completing additional, time-intensive departmental duties, such as teaching. My teaching portfolio boasts four years of teaching at different capacities (lectures, recitations, grading), at different levels (calculus 1, 2, 3, differential equations, linear algebra, and discrete math), for three terms per year (Spring, Summer, and Fall). My commitment to teaching resulted in a nomination for the Elizabeth Baranger teaching award, the most prestigious teaching award in the university.

As part of my application I include a research statement, curriculum vitae, and a list of publications. Please let me know if there is anything else I can provide, and I look forward to hearing from you.

\closing{Sincerely,\\
\fromsig{\includegraphics[scale=.075]{signature.png}}\\
Youngmin Park
}
%----------------------------------------------------------------------------------------
%	OPTIONAL FOOTNOTE
%----------------------------------------------------------------------------------------

% Uncomment the 4 lines below to print a footnote with custom text
%\def\thefootnote{}
%\def\footnoterule{\hrule}
%\footnotetext{\hspace*{\fill}{\footnotesize\em Footnote text}}
%\def\thefootnote{\arabic{footnote}}

%----------------------------------------------------------------------------------------

\end{letter}

\end{document}
